\subsection{Student residence 17.~listopadu}

The whole building complex of the student residence and refectory 17. listopadu
contains two dormitory buildings (buildings A and B) and the building of a
former refectory (building C, nowadays in possession of the Faculty of
Humanities and it is being reconstructed).
There are two basement floors under these buildings and you can find there the
refectory, technical support, bed linen deposit, gym, music rehearsal room and
lots of other more or less used space.

The basement of the buildings A and B is connected, although it is not easy to
find the way in the beginning.

\subsubsection{Building A}
The building A is the higher one (20 floors + ground floor) and is also closer
to the road. There are four lifts (but frequently just three of them are working
and sometimes, usually when you need it the most, even less) and two emergency
staircases.
There is an unwritten rule that says that you should not use a lift if you do
not need to go to the fourth floor and below, to go just one floor up
or down or to take it from the fifth and below floors to get to the
ground floor. You also should go to the -1 floor only during the lunch hours.
All of these trips are much faster if you use the stairs and it will save you
from the hateful stares of other passengers that usually go to the double-digit
floors.
But of course, the exceptions are allowed in the case of injury, exhaustion etc.

The windows either face the east (with a picturesque view of Magistrála) or the
west (with a view of the building C and tunel Blanka and the traffic jams
occurring there).
The east side is more popular with the residents as it is much cooler during the
summer.

The residents of the building A's twentieth floor about twice a year organize
parties in the corridors (although usually without permission of the
authorities). Do not rather move there if there is a chance this would bother
you.

There is a computer laboratory (\uv{lab}) on the ground floor. You will find
more information about it in the chapter about labs.
The offices are also located behind the lifts.


\subsubsection{Building B}
The building B has 16 floors + ground floor.

There is a dormitory council's room, piano room, gym and music rehearsal room in
the basement. On the ground floor you can find a study room, small shop for the
dormitory residents, ecumenical room and the office of the dormitory chief.

There are two small lifts and one emergency lift. It very often happens that at
least one of them does not go to all floors (especially the 15th floor) or does
not work at all. There is also an emergency staircase.
While you can frequently see a sign with \uv{MIMO PROVOZ} (out of order) on the
lifts in the building A, the lifts in the building B usually have the sign
\uv{VÝTAH V PROVOZU} (the lift is working). It is up to you to figure out the
hidden meaning of this.

The windows either face the east with the view of Milada (a former villa,
nowadays in ruins, and former squat that should have been demolished a long time
ago) or the west with the view of a hill.
Both of the sides of the building are hidden from the sun as the building A is
higher and offers shade.

\subsubsection{Former refectory (building C)}
Nowadays it is under construction. The Faculty of Humanities should move here
once the building is reconstructed because this faculty does not have its own
building.
Because all these buildings are interconnected, even the students that do not
have rooms close to the construction site are often disturbed by the noise.


\subsubsection{Magistrála}
Magistrála leads from Jižní Město to Prosek. It is a real blessing for our
dormitory as it makes a lot of noise and is like a poison for our lungs. Its
uses are endless.
Watching the road maintenance workers trying to get rid off the snow is a
distraction during the winter. Counting the passing cars can distract you while
studying for the exams. Plus it emits a nice low frequency sound helping you to
fall asleep.
Its illumination resembles the Greek letter epsilon.

A crowd of drivers, moving at the speed of a slow pedestrian, forms at
Magistrála every morning.
This might be beneficial for some people's confidence as they are easily able to
get ahead of the unnerved drivers.

After endless years of construction the tunnel complex Blanka has been connected
to Magistrála. In the afternoon there are usually traffic jams outside of it.


\subsubsection{Technical specialities}

Some people claim that if a fixed overcritical number of windows, all of them
above each other, is opened, the dormitory building falls down. It has not been
experimentally verified yet.
The building itself has the qualities of a chimney. While there used to be cold
on the lowest floors, you never needed to turn on your heating on the eighteenth
floor even if the weather was freezing, despite the fact that the building was
not thermally insulated at the time (residents staying on the highest floor used
to have different problems of course -- if you turned your heating on, you were
often unable to turn it off again).

The meteorologists will certainly gladly tell you why the wind is always strong
in the vicinity of the dormitory.
While there is just a gentle breeze blowing in Prague, the students walking
around the building A have problems to stay on their feet as the wind is very
strong. Only the luckiest ones manage to take a small step forward.
The whole trip along the (shorter!) wall of the building A might take a few
minutes. After a few windy days even the most stubborn student understands why
it is not always good to open the windows without checking whether you will not
need another three people to help you close them afterwards.

The construction works in the vicinity of the dormitory generally start at the
beginning of the examination period; it is hard to say whether this is just a
coincidence or an evil plan.
For example the anti-flood measures were built in the winter of 2009, the
reconstruction of the entrance hall took place in the summer of 2009, the
foundations of a building near the dormitory were laid in the summer of 2014,
the same went for the construction of the tunnel complex Blanka or the
reconstruction of the building C.


\subsubsection{Eduroam}
It is much easier to connect to Eduroam (university wifi network) now. You can
set you eduroam password in the CAS system. After this everything tends to work
the same way as if you tried to connect to any other wifi.
ATTENTION! your login name is: name@cuni.cz or isicnumber@cuni.cz.


\subsubsection{Dormitory council}
It is elected by the dormitory residents and represents their interests on the
regular meetings with the dormitory management.
You will find many useful information on their website -- important phone
numbers, forms, rules, lift etiquette, opening hours and a simulator for
organizing furniture.
Website address: \url{listopad.koleje.cuni.cz/lang/en}.

\subsubsection{Visits}
If your relatives or girlfriend/boyfriend are supposed to visit and you want to
accommodate them in your room, it is possible. You basically need to fulfil only
2 conditions. 1) Ask your roommate, if (s)he agrees and allows it and 2)
announce it in the accommodation office and at the reception -- or just at the
reception. The visitor will pay only 50 CZK per night.
Anyway, do not forget to keep the night silence (between 10 p.m. and 7 a.m.)
also during his/her visit.